\documentclass[11pt]{article}
\usepackage[margin=1in]{geometry}
\title{SNc, Intact Hemisphere --- Interpretation}
\author{GSE233866, 6-OHDA lesion cohort}
\date{}

\begin{document}
\maketitle

The intact-hemisphere SNc network, built from 8{,}637 cells (the largest
single-condition population in this project), independently selects a soft
power of 6 ($R^2 = 0.959$) --- higher than the healthy-baseline SNc network's
power of 5, consistent with this being a different animal cohort under a
different experimental arm rather than a repeat of the same data.

CACNA1D falls in the \textbf{blue} module here with $\mathrm{kME}=0.346$.
\emph{(Correction, 2026-08-13: this document previously compared against
0.048/0.094 for the two healthy-baseline networks and called 0.346
"substantially stronger than either" --- both those figures were misread
from the wrong \texttt{kME\_*} column; the corrected healthy-baseline values
are 0.442 (SNc) and 0.296 (VTA).)}

\textbf{This network's kME is not directly comparable to the
healthy-baseline networks' kME, and should not be read as a before/after
value.} The healthy-baseline arm (6 untreated animals) and this lesion arm
(a different 6 animals) are separate cohorts, each independently clustered
into its own network with its own module composition --- a numeric
difference between them reflects cohort/batch variation as much as
biology. The comparison this network \emph{is} built for is the paired one
within this same lesion-arm cohort: intact (this network, 0.346) versus
lesioned (\texttt{snc\_lesioned\_network/}, 0.398), same six animals,
contralateral hemispheres each. That comparison is the basis for the
"connectivity increases after lesioning" claim, and the formal version of
it is the DME test in \texttt{lesioned\_vs\_intact\_combined\_dme\_snc/}.
This network's main use on its own is as a sanity check that CACNA1D's
connectivity in intact tissue is not itself an outlier before interpreting
how it changes after lesioning.

\end{document}
